% #############################################################################
% This is Chapter 3 - Research Background
% !TEX root = ../main.tex
% #############################################################################
\fancychapter{Research Background}
\label{chap:background}

This chapter presents the theoretical background of the topics associated with the research. The topics covered are Artificial Intelligence, the Software Development Life Cycle, AI in Software Development, Vibe Coding, and Spec-Driven Development.


\section{Artificial Intelligence}

Artificial Intelligence (AI) refers to systems capable of performing tasks that traditionally require human cognitive abilities, such as reasoning, decision-making, pattern recognition, or language understanding. Early AI systems were rule-based, relying on explicitly defined logic; more recent developments are data-driven and statistical, with models improving their performance by identifying patterns in large datasets.

Generative Artificial Intelligence (GenAI) is the subset of AI designed to create new content rather than classify or retrieve existing information, producing text, images, audio or code by predicting likely outputs given prior context. That generative capacity is what positions GenAI as a significant innovation in domains relying on iteration, prototyping and creative exploration, including software development.

Large Language Models (LLMs) are a specific class of GenAI systems trained on vast corpora of text and source code. Built on transformer architectures, their primary function is next-token prediction, but that capability enables more complex tasks such as summarization, explanation, reasoning, conversational interaction, and code generation. Within software engineering, LLMs are relevant because they can interpret programming intent expressed in natural language and generate syntactically valid code in a variety of programming languages, as well as explain, refactor or restructure existing code and produce tests and documentation with far less manual effort. These capabilities position LLMs as collaborative tools supporting both novice and experienced developers across the stages of software development. \cite{Ahmed:2025ai,AlHashimi:2025er, Bai:2025ci, Banh:2025cf, Basha:2025tt, Chen:2025es, Dong:2024sc, Durrani:2024ad, Farhanna:2025eu, Gao:2025cc, Han:2025mg, Haque:2025la, Hemmat:2025rd, Hou:2025cl, J.Pries-Heje:2008SD, Jabrw:2025ss, Jensen:2025me, KarlovsKarlovskis:2024ga, Kim:2025ao, Kim:2025is, Li:2025sc, Moore:2025vc, Pendyala:2025pi, Qiu:2025ft, Ramos:2025aa, Terragni:2025fa, Tkachuk:2025dc, Wang:2025rs, Weyssow:2025ep, Yoo:2025sa, Zamfiroiu:2025ec, Zhang:2025ea}

\section{Software Development Life Cycle}

The Software Development Life Cycle (SDLC) describes the structured phases involved in building and maintaining software systems. Therefore, the SDLC is commonly structured into six main phases: requirements analysis, where system functionality is identified through requirements engineering; design, which defines the system architecture, data structures, and component interactions; implementation, focused on writing and integrating code; testing, where correctness, reliability, and performance are verified; deployment, which delivers the system to end users; and maintenance, involving ongoing updates, improvements, and operational support throughout the system’s lifetime.

The SDLC is central to this research because LLMs can engage differently across its stages: they can support implementation by generating code, testing by producing test cases, and maintenance by summarizing or refactoring existing code.

Software development processes have evolved to emphasize iterative delivery and continuous integration, supported by methodologies such as Agile and DevOps. Despite these improvements, much of the work remains manual, repetitive and time-intensive: debugging, documenting code, maintaining consistency across repositories, and handling routine implementation. As systems scale in complexity, the pressure to increase productivity without compromising maintainability or security grows with them, and that pressure is a key factor driving interest in AI-assisted development workflows. \cite{Zhang:2025ea, Ahmed:2025ai, Jabrw:2025ss, Terragni:2025fa}

\section{AI in Software Development}

Advances in LLMs have introduced new forms of interaction between developers and software tooling, shifting the developer from writing code directly to supervising and guiding automated generation. Studies indicate that LLMs can accelerate prototyping, clarify unfamiliar codebases, and help maintain documentation quality, which is especially valuable in early-stage development or when exploring new frameworks and libraries.

The integration is increasingly characterized as collaborative rather than as full automation: LLMs support pattern recognition, summarization and structural reasoning, while the developer remains responsible for architectural decisions, validation, integration, and alignment with project requirements. Research in this area focuses on how workflows adapt to that shift, how developers form trust in AI outputs, and how responsibilities are distributed between human judgment and automated generation. These themes are further examined in \Cref{chap:slr}. \cite{Ahmed:2025ai, Zhang:2025ea, Durrani:2024ad, Jabrw:2025ss, Banh:2025cf, Farhanna:2025eu, KarlovsKarlovskis:2024ga, Terragni:2025fa}

\section{Vibe Coding}

Vibe Coding refers to a style of software development in which developers describe their intentions, requirements, or constraints in natural language, and a generative model produces corresponding code or related artifacts. The developer does not begin by writing syntactically precise code: development proceeds conversationally, the model proposing implementations that the developer then evaluates for correctness, aligns with project requirements, and refines by modifying either the prompt or the produced code.

Vibe Coding therefore shifts part of the programming effort from manual construction to specifying intent and assessing alternatives. It does not replace traditional programming but reframes how initial structure and exploratory code are produced, with implications for developer workflows, cognitive load, and the distribution of effort across the SDLC that are examined in \Cref{chap:slr}. \cite{Moore:2025vc, Zamfiroiu:2025ec}

\section{Spec-Driven Development}
\label{sec:background_sdd}

Spec-driven development is an emerging engineering practice in which a persistent, versioned specification, rather than the source code itself, is treated as the primary artefact that both developers and AI agents work against. Instead of prompting a model directly against a codebase and accepting whatever it produces, the practice interposes a written specification, of requirements, design decisions, and constraints, that the model must consult and satisfy before code is generated, and that is kept in sync with the codebase as it evolves. Industry examples of this pattern include AWS Kiro's three-stage requirements, design, and tasks workflow, GitHub's Spec-Kit, which adds a "constitution" of immutable project principles that generated code may not violate, and the Tessl Framework, in which, as its own materials put it, the specification itself becomes the maintained artefact rather than the code \cite{Thoughtworks:2025SD}.

The practice is tracked by Thoughtworks' Technology Radar as an emerging AI-assisted engineering technique, named specifically "spec-driven development" \cite{Thoughtworks:2025SD}, and has begun to receive academic attention: Farrag frames it as specification-driven governance, arguing that it addresses what is termed the productivity-reliability paradox in AI-augmented development, though this framing is currently reported only in a preprint that has not yet completed peer review \cite{Farrag:2026sd}. The practice responds directly to the risk profile characterised in \Cref{sec:motivation}: narrowing and anchoring the context a model is permitted to act upon is a governance mechanism, not merely a convenience, since it gives generated output a stable reference to be validated against rather than leaving correctness to be inferred from the prompt alone. This precedent is the direct grounding for Spec-Anchored Continuity, the persistent specification mechanism at the centre of the framework proposed in \Cref{chap:proposal}.
